% Page 051

\seite{51}

\abschnitt{Spaziergang}
\pstart
\margmark{Am 3 Septe}%
mber machte hiesige Schule einen Spaziergang nach St.\ Thomas

\pend
\abschnitt{Geschehen 8. 1. 14}
\pstart


\pend
\abschnitt{Kirchlichkeiten}
\pstart

\margmark{Der Geburt\\feiert von\\Kindern wu\\und die Fe\\jener Kron}%
stag Sr. Majestät des Kaisers wurde im Beisein ge\ohb
 den 2 Vätern derselben im Saale des Gasthauses. Den\ob
rde{[n]} Geschenke\unsicher{} es wurden vaterländische Lieder gesungen\ob
ste vorgetragen. Die Aufsätze an der Kinder galt\ob
e von Kuchen.

\pend
\abschnitt{Osterprüfung}
\pstart

\margmark{wurde am 8\\mit Segen}%
. März durch den Ortsschulinspek{[tor]} gegen Abend\unsicher\ob
vorgehalten.

\pend
\jahresueberschrift{Das Schuljahr 1914/15}
\pstart

\margmark{Aufgelösen\\2 Mädchen.\\die Genomm\\Die Schüle\\Alle Kinde}%
 wurden in hiesigen Schule 4 Kinder 2 Knaben\ob
 Darüber mit denen für die Landwirtschaft\ob
en haben an 9 Kinder 2 stärkere Klasse.\ob
rzahl beträgt 47, 11 Knaben, 11 Mädchen.\ob
r sind katholisch.


Besondere Ereignisse\ob
\margmark{Die zweite\\Krieg brac\\mahlen zu\\Slaverei v\\leider Öst\\England, F\\die gerne\\es war auß\\und suchte\\Volk gut d\\Volk mit A}%
 Hälfte des Jahres 1914 brachte Deutschland große Nöte, der\ob
h auf das Volk es galten den Franzosen und dessen\unsicher{} Ge\ohb
ermahnet. Da die Reden die Bedingungen nicht erfüllten, die\ob
on ihnen forderte, daß des König in Deutschland, als Vater\ob
erreichs, nur bereit, zu stehen. Die Stimmen der Russen,\ob
rankreich hatten Längs einen Kniff und Deutschlander\ob
in aller gebildeter Weise dem Lande den Frieden erhalten\ob
en eifrig. Den erlosch ja, das Mounts zu großen hat\ob
 einst Kaiser und Deutschland, als Rußland bereit das\ob
er Waffen gerufen hatte. Es waren gen Tage das\ob
ngst, alle immer sichter nun, es warte dem Kaiser
\pend
